\documentclass[12pt,a4paper]{article}
\usepackage{pdfpages}
\usepackage{float}
\usepackage[utf8]{inputenc}
\usepackage[T1]{fontenc}
\usepackage[french]{babel}
\usepackage{geometry}
\usepackage{fancyhdr}
\usepackage{lmodern}
\usepackage[most]{tcolorbox}
\usepackage{xcolor}
\usepackage{graphicx}
\usepackage[hidelinks]{hyperref}
\geometry{
    top=2cm,
    bottom=2cm,
    left=2cm,
    right=2cm
}

% Header
\pagestyle{fancy}
\fancyhf{}

\fancyhead[L]{MERN Stack}
\fancyfoot[C]{\thepage}
\fancyhead[R]{étudiante : Syrine Zmerli}
\renewcommand{\headrulewidth}{0.4pt}
\renewcommand{\footrulewidth}{0pt}

\begin{document}
\pagenumbering{roman}
\includepdf[pages=-]{page de garde}
\tableofcontents %fr
\clearpage
\listoffigures  %fr
\clearpage

\pagenumbering{arabic}
\begin{center}

    {\LARGE \textbf{Compte Rendu MERN Stack -- Séance 4}}\\[0.2cm]

    {\large \textbf{Sécurisation de l'API REST : Authentification JWT, Hachage Bcrypt \& Middlewares}}\\[0.15cm]

    \vspace{0.3cm}

    \hrule

\end{center}

\vspace{0.5cm}

\section{La Prédiction}

\begin{tcolorbox}[
    colback=orange!5,
    colframe=orange!80!black,
    boxrule=1pt,
    arc=5pt,
    left=8pt,
    right=8pt,
    top=5pt,
    bottom=5pt
]
\textcolor{orange!80!black}{\textbf{Prédisez.}}
\quad Pourquoi est-il strictement interdit d'utiliser un algorithme de chiffrement réversible (comme AES) ou une simple fonction de hachage rapide (comme MD5 ou SHA-256 brute) pour enregistrer les mots de passe ? Que se passe-t-il si un attaquant utilise des tables arc-en-ciel (rainbow tables) précalculées ?
\end{tcolorbox}

\textbf{--> Réponse \& Analyse :} \\
Il est strictement interdit d'utiliser ces méthodes pour les raisons suivantes :
\begin{itemize}
    \item \textbf{Le chiffrement réversible (AES) :} Le chiffrement symétrique comme AES est à double sens : avec la clé secrète, on peut retrouver le mot de passe original en clair. Si un attaquant parvient à compromettre le serveur, accéder au code source ou extraire la variable d'environnement contenant la clé secrète, il pourra déchiffrer tous les mots de passe de tous les utilisateurs immédiatement.
    \item \textbf{Les fonctions de hachage rapides (MD5, SHA-256) :} Ces algorithmes ont été conçus pour vérifier rapidement l'intégrité de fichiers et sont optimisés pour la vitesse. Un ordinateur avec une carte graphique moderne (GPU) peut calculer des milliards de condensats SHA-256 par seconde. Un attaquant peut donc retrouver un mot de passe très facilement par attaque par force brute ou par dictionnaire.
    \item \textbf{Les tables arc-en-ciel (Rainbow Tables) :} Sans sel (salt), deux mots de passe identiques produisent exactement le même hash. Une table arc-en-ciel contient des millions de correspondances précalculées entre mots de passe courants et leurs empreintes de hachage. Si un attaquant vole notre base de données sans sel, il n'a même pas besoin de calculer : il fait une recherche instantanée dans sa table pour retrouver le mot de passe en clair.
    \item \textbf{La solution avec Bcrypt :} \texttt{bcryptjs} résout ces failles en combinant un \textbf{sel aléatoire (salt)} unique pour chaque utilisateur (ce qui rend les tables arc-en-ciel totalement inutiles), et un \textbf{facteur de coût (work factor)} qui ralentit volontairement le calcul pour bloquer la force brute.
\end{itemize}

\clearpage

\section{Campagne de Tests Méthodique avec Postman (Étape 8)}

\subsection{Tableau Récapitulatif des Tests}

\renewcommand{\arraystretch}{1.2}
\begin{table}[H]
\centering
\begin{tabular}{|l|l|p{6.5cm}|l|}
\hline
\textbf{Méthode} & \textbf{Endpoint} & \textbf{Conditions / En-têtes} & \textbf{Statut attendu} \\
\hline
POST & \texttt{/api/users/register} & Body avec nom, email, password ($\ge 6$) & 201 Created \\
\hline
POST & \texttt{/api/users/login} & Mêmes identifiants valides $\rightarrow$ Récupération du JWT & 200 OK \\
\hline
GET & \texttt{/api/users/profile} & Sans en-tête Authorization & 401 Unauthorized \\
\hline
GET & \texttt{/api/users/profile} & Header : \texttt{Authorization: Bearer <token>} & 200 OK \\
\hline
POST & \texttt{/api/articles} & Header avec Token + Body \{ title, content \} & 201 Created \\
\hline
GET & \texttt{/api/users} & Token avec role: "user" $\rightarrow$ Test RBAC admin & 403 Forbidden \\
\hline
\end{tabular}
\caption{Matrice des tests d'authentification et de contrôle d'accès}
\end{table}

\subsection{Vérifications Postman}

\textbf{1. Inscription d'un nouvel utilisateur (POST /api/users/register)}
\begin{figure}[H]
    \centering
    \includegraphics[width=0.85\linewidth]{assets/postman_register.png}
    \caption{POST /api/users/register - Statut 201 Created avec token généré}
\end{figure}
\textbf{Explication :} \\
--> L'utilisateur s'enregistre avec son nom, email, mot de passe et age. Le mot de passe est haché avec un sel avant d'être sauvegardé dans MongoDB. L'API renvoie les données de l'utilisateur avec son rôle ("user") et son token JWT.

\vspace{0.3cm}
\textbf{2. Connexion et récupération du Token (POST /api/users/login)}
\begin{figure}[H]
    \centering
    \includegraphics[width=0.85\linewidth]{assets/postman_login.png}
    \caption{POST /api/users/login - Statut 200 OK avec récupération du token}
\end{figure}
\textbf{Explication :} \\
--> Le contrôleur recherche l'utilisateur par son email, puis compare le mot de passe reçu avec le hash stocké en base grâce à \texttt{bcrypt.compare}. Si la comparaison réussit, le serveur renvoie le token JWT.

\clearpage

\textbf{3. Rejet de l'accès sans Token (GET /api/users/profile)}
\begin{figure}[H]
    \centering
    \includegraphics[width=0.85\linewidth]{assets/postman_profile_no_token.png}
    \caption{GET /api/users/profile sans token - Statut 401 Unauthorized}
\end{figure}
\textbf{Explication :} \\
--> Le middleware \texttt{protect} vérifie la présence du token dans l'en-tête \texttt{Authorization}. Comme aucun token n'a été fourni, la requête est rejetée avec le message \texttt{"Non autorise , aucun token fourni"}.

\vspace{0.3cm}
\textbf{4. Accès au profil avec Token (GET /api/users/profile)}
\begin{figure}[H]
    \centering
    \includegraphics[width=0.85\linewidth]{assets/postman_profile_with_token.png}
    \caption{GET /api/users/profile avec Bearer Token - Statut 200 OK}
\end{figure}
\textbf{Explication :} \\
--> En ajoutant le header \texttt{Authorization: Bearer <token>}, le middleware vérifie la signature avec \texttt{JWT\_SECRET}, récupère l'utilisateur en base sans son mot de passe (\texttt{.select('-password')}) et l'injecte dans \texttt{req.user}. Le contrôleur renvoie alors le profil connecté.

\clearpage

\textbf{5. Création protégée d'un article (POST /api/articles)}
\begin{figure}[H]
    \centering
    \includegraphics[width=0.85\linewidth]{assets/postman_create_article.png}
    \caption{POST /api/articles - Statut 201 Created (auteur lié automatiquement)}
\end{figure}
\textbf{Explication :} \\
--> L'auteur de l'article n'est plus envoyé dans le body par le client : il est extrait de manière sécurisée depuis le token validé via \texttt{req.user.\_id}. Ceci empêche toute usurpation d'identité.

\vspace{0.3cm}
\textbf{6. Rejet d'un utilisateur simple sur une route admin (GET /api/users)}
\begin{figure}[H]
    \centering
    \includegraphics[width=0.85\linewidth]{assets/postman_admin_forbidden.png}
    \caption{GET /api/users avec rôle user - Statut 403 Forbidden}
\end{figure}
\textbf{Explication :} \\
--> Le middleware \texttt{adminOnly} vérifie que \texttt{req.user.role === 'admin'}. Comme l'utilisateur a le rôle simple "user", l'accès est bloqué avec une erreur 403 et le message \texttt{"Acces refuse : droits administrateur requis"}.

\clearpage

\subsection{Point de Contrôle MongoDB Atlas}
\begin{figure}[H]
    \centering
    \includegraphics[width=0.9\linewidth]{assets/atlas_hash.png}
    \caption{Document MongoDB Atlas avec le mot de passe haché (\$2a\$10\$...)}
\end{figure}
\textbf{Explication :} \\
--> En ouvrant l'interface MongoDB Atlas dans la collection \texttt{users}, on constate que le mot de passe en clair n'apparaît nulle part. Il est stocké sous forme d'une empreinte de 60 caractères commençant par \texttt{\$2a\$10\$}, ce qui confirme le bon fonctionnement du salage et du hachage avec Bcrypt.

\clearpage

\section{Exercice 1 : Mise à Jour Sécurisée du Mot de Passe}

\begin{tcolorbox}[
    colback=green!5,
    colframe=green!60!black,
    boxrule=1pt,
    arc=5pt,
    left=8pt,
    right=8pt,
    top=5pt,
    bottom=5pt
]
\textbf{Exercice 1 : Mise à Jour Sécurisée du Mot de Passe}
\begin{enumerate}
    \item Déclarez la route \texttt{PUT /api/users/update-password} protégée par le middleware \texttt{protect}.
    \item Le corps de la requête contient deux champs : \texttt{oldPassword} et \texttt{newPassword}.
    \item Vérifiez d'abord que le mot de passe actuel correspond bien à celui stocké en base en faisant appel à \texttt{bcrypt.compare(oldPassword, user.password)}. En cas d'échec, renvoyez une erreur 401 Unauthorized avec le message : \texttt{"Ancien mot de passe incorrect"}.
    \item Si l'ancien mot de passe est vérifié, générez un nouveau sel, hachez \texttt{newPassword}, sauvegardez l'utilisateur avec \texttt{await user.save()} et renvoyez un code 200 OK avec un message de succès.
    \item Validez le bon fonctionnement avec Postman en vous reconnectant avec le nouveau mot de passe.
\end{enumerate}
\end{tcolorbox}

\subsection{Code Implémenté}
\begin{figure}[H]
    \centering
    \includegraphics[width=0.95\linewidth]{assets/code_ex1.png}
    \caption{Code de la fonction updatePassword dans userController.js}
\end{figure}

\subsection{Vérification Postman}
\begin{figure}[H]
    \centering
    \includegraphics[width=0.85\linewidth]{assets/postman_ex1_wrong.png}
    \caption{Test avec ancien mot de passe incorrect - Statut 401 Unauthorized}
\end{figure}

\begin{figure}[H]
    \centering
    \includegraphics[width=0.85\linewidth]{assets/postman_ex1_success.png}
    \caption{Test avec mot de passe valide - Statut 200 OK}
\end{figure}

\textbf{Explication :} \\
--> J'ai créé la fonction \texttt{updatePassword} dans \texttt{userController.js}. Comme \texttt{req.user} extrait l'utilisateur sans le mot de passe (\texttt{-password}), je recherche à nouveau l'utilisateur par son id avec \texttt{User.findById(req.user.\_id)} pour obtenir son mot de passe haché. Ensuite j'utilise \texttt{bcrypt.compare(oldPassword, user.password)} : si c'est faux je renvoie une 401, sinon je génère un nouveau sel, je hache le nouveau mot de passe, je fais \texttt{await user.save()} et je renvoie un code 200.

\clearpage

\section{Exercice 2 : Route Filtrée des Articles Personnels \& Rôles}

\begin{tcolorbox}[
    colback=green!5,
    colframe=green!60!black,
    boxrule=1pt,
    arc=5pt,
    left=8pt,
    right=8pt,
    top=5pt,
    bottom=5pt
]
\textbf{Exercice 2 : Route Filtrée des Articles Personnels \& Attribution de Rôles}
\begin{enumerate}
    \item Créez une nouvelle route protégée \texttt{GET /api/articles/my-articles} placée avant la route \texttt{GET /api/articles/:id} dans \texttt{articleRoutes.js}.
    \item Dans le contrôleur associé \texttt{getMyArticles}, filtrez les articles appartenant exclusivement à l'utilisateur connecté grâce à la requête Mongoose :
    \texttt{const articles = await Article.find(\{ author: req.user.\_id \});}
    \item Protection contre l'escalade de privilèges : Dans la fonction \texttt{registerUser}, faites en sorte que le champ \texttt{role} ne puisse jamais être défini à 'admin' par un appel d'inscription public (forcez toujours \texttt{role = 'user'}).
    \item Créez une route d'administration dédiée \texttt{PUT /api/users/:id/role} protégée par \texttt{protect} et \texttt{adminOnly} permettant exclusivement à un administrateur d'élever un utilisateur au rang d'admin.
\end{enumerate}
\end{tcolorbox}

\subsection{Code Implémenté}
\begin{figure}[H]
    \centering
    \includegraphics[width=0.95\linewidth]{assets/code_ex2_routes.png}
    \caption{Positionnement de /my-articles avant /:id dans articleRoutes.js}
\end{figure}

\begin{figure}[H]
    \centering
    \includegraphics[width=0.95\linewidth]{assets/code_ex2_controller.png}
    \caption{Fonction getMyArticles et fonction updateUserRole}
\end{figure}

\clearpage

\subsection{Vérification Postman}

\begin{figure}[H]
    \centering
    \includegraphics[width=0.85\linewidth]{assets/postman_my_articles.png}
    \caption{GET /api/articles/my-articles - Récupération exclusive des articles personnels}
\end{figure}

\begin{figure}[H]
    \centering
    \includegraphics[width=0.85\linewidth]{assets/postman_update_role.png}
    \caption{PUT /api/users/:id/role - Élévation de rôle par un administrateur}
\end{figure}

\textbf{Explication :} \\
--> Pour la question 1 et 2, il est essentiel de placer la route \texttt{/my-articles} avant \texttt{/:id}. Si on la mettait après, Express interpréterait la chaîne \texttt{"my-articles"} comme un paramètre \texttt{id} de type ObjectId, ce qui déclencherait une erreur Mongoose \texttt{CastError}. Dans \texttt{getMyArticles}, on filtre par \texttt{\{ author: req.user.\_id \}} pour que chaque utilisateur n'ait accès qu'à ses propres articles.\\
--> Pour la question 3, dans \texttt{registerUser}, j'ai fixé en dur \texttt{role: 'user'} dans la création du document pour empêcher qu'un utilisateur malveillant s'enregistre comme administrateur en envoyant \texttt{"role": "admin"} dans le JSON.\\
--> Pour la question 4, la route \texttt{PUT /api/users/:id/role} est protégée par \texttt{protect} et \texttt{adminOnly}. Seul un administrateur connecté peut changer le rôle d'un compte.

\clearpage

\section{Exercice 3 : Gestion de l'Expiration du Jeton (Token Expired)}

\begin{tcolorbox}[
    colback=green!5,
    colframe=green!60!black,
    boxrule=1pt,
    arc=5pt,
    left=8pt,
    right=8pt,
    top=5pt,
    bottom=5pt
]
\textbf{Exercice 3 : Gestion Explicite de l'Expiration du Jeton (Token Expired)}
\begin{enumerate}
    \item Dans votre middleware \texttt{protect} (\texttt{authMiddleware.js}), capturez l'erreur dans le bloc \texttt{catch(err)}.
    \item Si \texttt{err.name === 'TokenExpiredError'}, renvoyez un statut 401 Unauthorized avec le message JSON précis :
    \texttt{\{ "message": "Votre session a expire , veuillez vous reconnecter" \}}
    \item Pour tester facilement, modifiez temporairement l'option d'expiration dans \texttt{generateToken} à '10s'. Connectez-vous, attendez 15 secondes, puis effectuez une requête sur \texttt{/api/users/profile}. Vérifiez la capture de l'expiration. Rétablissez ensuite à '1d'.
\end{enumerate}
\end{tcolorbox}

\subsection{Code Implémenté}
\begin{figure}[H]
    \centering
    \includegraphics[width=0.95\linewidth]{assets/code_ex3.png}
    \caption{Capture de TokenExpiredError dans authMiddleware.js}
\end{figure}

\subsection{Vérification Postman}
\begin{figure}[H]
    \centering
    \includegraphics[width=0.85\linewidth]{assets/postman_token_expired.png}
    \caption{Test après 15 secondes d'attente - Statut 401 et message d'expiration}
\end{figure}

\textbf{Explication :} \\
--> Lorsque la durée de validité du token dépasse la valeur définie par \texttt{expiresIn}, la méthode \texttt{jwt.verify()} lève automatiquement l'exception \texttt{TokenExpiredError}. Dans le bloc \texttt{catch}, j'ai ajouté une condition sur \texttt{err.name === 'TokenExpiredError'} pour renvoyer le message précis \texttt{"Votre session a expire , veuillez vous reconnecter"}. Après vérification du test avec 10s, j'ai rétabli la durée normale à \texttt{'1d'}.

\clearpage

\section{Lien Github}
\href{https://github.com/SyrineZmerli/Mern_TP1/tree/master}{Lien Github vers mon repositoire : https://github.com/SyrineZmerli/Mern\_TP1}

\end{document}
